\documentclass[11pt]{article}
\usepackage{mathunal-base}
\mathunalfuente{serif}

\begin{document}

{\large\bfseries Parcial Final de Matemáticas Básicas --- selección múltiple}

\medskip
\noindent Escuela de Matemáticas. Universidad Nacional de Colombia, Sede Medellín.

\medskip
\noindent\textit{Examen de selección múltiple con única respuesta. Marque una sola opción por pregunta.}

\begin{enumerate}
\item Solo una de las siguientes proposiciones es verdadera. Indique cuál.
\begin{opciones}
\item $\sqrt{\sqrt{16}} = \pm 2$
\item $-\sqrt[4]{(-2)^4} = -2$
\item $\sqrt[4]{(-2)^4} = -2$
\item $\sqrt[n]{x^m} = x^{m/n}$ para todo $x \in \mathbb{R}$
\end{opciones}

\item El polinomio cuadrático $P(x) = 3x^2 + 2x - 8$ tiene:
\begin{opciones}
\item Dos ceros reales distintos: $\tfrac{4}{3}$ y $-2$
\item Un cero real $-2$
\item Dos ceros reales distintos: $-\tfrac{4}{3}$ y $2$
\item No tiene ceros reales
\end{opciones}

\item Considere la función definida por $f(x) = \dfrac{x + 3}{\sqrt{x^2 + x - 2}}$; su dominio es:
\begin{opciones}
\item $(-\infty, -2) \cup (1, \infty)$
\item $(-\infty, -3) \cup (-3, -1)$
\item $(-\infty, -3) \cup (1, \infty)$
\item $(-\infty, -9) \cup (-9, -1)$
\end{opciones}

\item Todos los siguientes son trinomios cuadrados perfectos, excepto uno; indíquelo.
\begin{opciones}
\item $\tfrac{9}{4}a^6 + 3a^3 - 1$
\item $1 + 9a^2 - 6a$
\item $4m^2 + 9n^4 - 12mn^2$
\item $25x^4 y^6 - 40x^2 y^3 z^4 + 16z^8$
\end{opciones}

\item Luego de simplificar la expresión $\sqrt[4]{5\sqrt{25}}$, se obtiene:
\begin{opciones}
\item $\sqrt[3]{5}$
\item $\sqrt[3]{125}$
\item $\sqrt5$
\item $\sqrt{25}$
\end{opciones}

\item Si eliminamos los radicales del denominador de la expresión $\dfrac{h}{\sqrt{x + h} - \sqrt{x}}$ obtenemos:
\begin{opciones}
\item $\sqrt{x + h} - \sqrt{x}$
\item $\sqrt{x + h} + \sqrt{x}$
\item $\dfrac{\sqrt{x}}{x}$
\item $\dfrac{\sqrt{x + h}}{2}$
\end{opciones}

\item Después de eliminar los denominadores de la ecuación
\[
\frac{x}{3} + \frac{2x - 1}{5} = 6 - \frac{3x + 2}{10}
\]
queda:
\begin{opciones}
\item $10x + 6(2x - 1) = 180 - 3(3x + 2)$
\item $5x + 3(2x - 1) = 60 - (3x + 2)$
\item $10x - 6(2x - 1) = 180 + 3(3x + 2)$
\item $5x - 3(2x - 1) = 60 + (3x + 2)$
\end{opciones}

\item Un recipiente en forma de cono circular recto, con vértice en la base, tiene un volumen total de $10\pi$ cm$^3$. Cuando este recipiente se llena de vino hasta un volumen igual a $\tfrac{2}{3}$ de su volumen total, la profundidad es $4$ cm. Entonces el radio del cono formado por el contenido de vino en el recipiente es:
\begin{opciones}
\item $2\sqrt5$ cm
\item $5$ cm
\item $4$ cm
\item $\sqrt{\tfrac{5}{3}}$ cm
\end{opciones}

\item Dos números reales suman $21$. Los números para los cuales el producto es $110{,}25$ son:
\begin{opciones}
\item $10{,}5$ y $10{,}5$
\item $10$ y $11$
\item $1$ y $20$
\item $18$ y $3$
\end{opciones}

\item En la notación de intervalos, el conjunto dado por $P = \{x \mid x \geq -2 \text{ y } x \leq 4\}$ se escribe:
\begin{opciones}
\item $(-\infty, -2] \cup [4, \infty)$
\item $[-2, 4]$
\item $[4, -2]$
\item $[-2, \infty)$
\end{opciones}

\item Todas las siguientes proposiciones son falsas, excepto una; indique cuál.
\begin{opciones}
\item Si $x$ es cualquier número real, entonces $x + x > x$
\item Si $x$ es cualquier número real, entonces $x^2 > x$
\item Si $x > 0$ y $y > 0$, entonces $x + y > y$
\item Si $y$ es cualquier número real diferente de $0$ y $\tfrac{4}{y} > 1$, entonces $4 > y$
\end{opciones}

\item La solución de la inecuación $|x + 2| > -x^2 - 4$ es:
\begin{opciones}
\item $-2 < x < 6$
\item $-6 < x < 2$
\item $\emptyset$
\item $\mathbb{R}$
\end{opciones}

\item Geométricamente, $\sqrt{17}$ es la longitud de la hipotenusa de un triángulo rectángulo cuyos catetos miden:
\begin{opciones}
\item $\sqrt5$ y $3$
\item $12$ y $5$
\item $1$ y $1$
\item $\sqrt{11}$ y $6$
\end{opciones}

\item La ecuación de un lado de un paralelogramo es $y = 2x - 3$. La pendiente del lado opuesto es:
\begin{opciones}
\item $3$
\item $-3$
\item $2$
\item $-2$
\end{opciones}

\item A un empleado de cierto almacén le pagan un salario mensual fijo de $200.000$, y $2.000$ adicionales por cada artículo que venda en el mes. Para obtener un salario superior a $600.000$ mensuales, deberá vender:
\begin{opciones}
\item A lo más $200$ artículos
\item Menos de $200$ artículos
\item Menos de $201$ artículos
\item Por lo menos $201$ artículos
\end{opciones}

\item Sea $g(x) = \dfrac{1 - 3x}{1 - x}$; su inversa $g^{-1}$ está dada por:
\begin{opciones}
\item $g^{-1}(x) = \dfrac{3 - x}{1 - x}$
\item $g^{-1}(x) = \dfrac{1 - x}{3 - x}$
\item $g^{-1}(x) = \dfrac{x - 3}{x - 1}$
\item $g^{-1}(x) = \dfrac{1 - 3x}{x - 1}$
\end{opciones}

\item El valor en radianes de $-1440^\circ$ es:
\begin{opciones}
\item $-8\pi$
\item $-16\pi$
\item $-4\pi$
\item $-24\pi$
\end{opciones}

\item La solución de la ecuación $\operatorname{sen} x \cos x = \tfrac{1}{2}$, para $0^\circ \leq x \leq 360^\circ$, es:
\begin{opciones}
\item $\{45^\circ, 225^\circ\}$
\item $\{45^\circ, 135^\circ\}$
\item $\{45^\circ, -45^\circ\}$
\item No tiene solución
\end{opciones}

\item La expresión $\operatorname{sen}^4\theta + 2\operatorname{sen}^2\theta \cos^2\theta + \cos^4\theta$ es idéntica a:
\begin{opciones}
\item $\operatorname{sen}^4\theta - \cos^4\theta$
\item $1$
\item $-1$
\item $2\operatorname{sen}^2\theta \cos^2\theta$
\end{opciones}

\item Para cualquier valor de $\theta$ se cumple que $\operatorname{sen}\theta$ es igual a:
\begin{opciones}
\item $1 - \cos\theta$
\item $\operatorname{sen}(\pi + \theta)$
\item $\operatorname{sen}(2\pi - \theta)$
\item $\cos\!\left(\tfrac{\pi}{2} - \theta\right)$
\end{opciones}
\end{enumerate}

\vfill
\noindent{\small\itshape Transcripción digital --- MathUNAL}

\end{document}
